\subsection{仿射 Springer 纤维的几何性质}

\subsubsection{非约化性} 若 $G$ 非平凡、$\gamma\in\frg(F)$ 正则半单，则归纳概形 $\tX_{\gamma}$ {\em 从不}约化。例如在 \S\ref{sss:reduction rs} 考虑的情形中，$\tX_{\gamma}$ 同构于 $\Gr_{T}=LT/L^{+}T$。在 \S\ref{sss:centralizer} 中已经看到，对非约化的环 $R$，$LT(R)=R\lr{t}^{\times}$ 中的元素可以有幂零的首项系数。因此 $\Gr_{T}(R)$ 并不只是 $\xcoch(T)$，后者才是 $\Gr^{\red}_{T}$。这说明 $\Gr_{T}$ 非约化，从而 $\tX_{\gamma}$ 非约化。

下面这个定理是关于 $\sX_{\gamma}$ 的基本有限性陈述。

\begin{theorem}[实质上是 Kazhdan 与 Lusztig {\cite[命题 2.1]{KL}}]\label{th:ft quot} 设 $\gamma\in \frg(F)$ 是正则半单元。则
\begin{enumerate}
\item 存在闭子概形 $Z\subset \sX_{\gamma}$，它在 $k$ 上射影，且 $\sX_{\gamma}=\cup_{\ell\in\L_{\gamma}}\ell\cdot Z$。
\item 归纳概形 $\sX_{\gamma}$ 是 $k$ 上局部有限型的概形。
\item $\L_{\gamma}$ 在 $\sX_{\gamma}$ 上的作用自由，且商 $\L_{\gamma}\bs\sX_{\gamma}$（作为 $k$ 代数上的 fppf 层）可由 $k$ 上真代数空间表示。
\end{enumerate}
\end{theorem}

下面分三步概述该定理的证明。

\subsubsection{第一步归约} 我们说明定理的第 (1) 部分蕴含 (2) 与 (3)。设 $Z$ 是如 (1) 中的射影子概形。为证 (2)，只需说明任一点 $x\in \sX_{\gamma}(k)$ 都有有限型的开邻域。借助 $\L_{\gamma}$-作用，可以假定 $x\in Z(k)$。由于 $Z$ 有限型，集合 $\Sigma:=\{\ell\in \L_{\gamma}|Z\cap\ell\cdot Z\neq\varnothing\}$ 有限。令 $U=\sX_{\gamma}-\cup_{\ell\notin\Sigma}\ell\cdot Z$，则 $U$ 是 $Z$ 的开邻域，从而也是 $x$ 的开邻域。此外 $U$ 包含在有限并 $\cup_{\ell\in\Sigma}\ell\cdot Z$ 中，从而包含在某个 Schubert 簇 $\Gr_{G,\leq\l}$ 中。因此 $U$ 是射影概形 $\sX_{\gamma}\cap\Gr_{G,\leq\l}$ 的开子集，故 $U$ 本身是有限型概形。为证 (3)，注意 fppf 层商 $\L_{\gamma}\bs\sX_{\gamma}$ 是分离代数空间，因为它是 $\sX_{\gamma}$ 对 étale 等价关系 $\L_{\gamma}\times\sX_{\gamma}\subset \sX_{\gamma}\times\sX_{\gamma}$（由作用映射与投影映射给出）的商。由 (1)，有射影概形 $Z$ 到 $\L_{\gamma}\bs\sX_{\gamma}$ 的满射，故 $\L_{\gamma}\bs\sX_{\gamma}$ 真。

\subsubsection{$\gamma$ 落在分裂环面中时 (1) 的证明}\label{sss:split case} 先考虑 $\gamma$ 落在分裂环面中的情形。经 $G(F)$-共轭，可假定 $\gamma\in\frt(F)$。此时 $\L_{\gamma}\cong\xcoch(T)$。固定含 $T$ 的 Borel 子群 $B$，令 $N$ 为 $B$ 的幂幺根。$LG$ 的 Iwasawa 分解给出
\begin{equation*}
\Gr_{G}=LN\cdot \L_{\gamma}L^{+}G/L^{+}G=\sqcup_{\l\in\xcoch(T)}LN\cdot \l(t)L^{+}G/L^{+}G.
\end{equation*}
令 $X:=\sX_{\gamma}\cap (LN\cdot L^{+}G/L^{+}G)\subset \sX_{\gamma}$。只需说明 $X$ 落在某个仿射 Schubert 簇 $\Gr_{\leq \l}$ 中，因为此时它在 $\Gr_{\leq\l}$ 中的闭包 $Z:=\ov{X}$ 便满足 (1) 的条件。为后文备用，注意平移 $\ell\cdot X$（$\ell\in \L_{\gamma}$）互不相交，且覆盖 $\sX_{\gamma}$。

固定 $T$ 关于 $B$ 的正根的一个次序，则可把元素 $u\in N(F)$ 唯一地写成
\begin{equation}\label{root gp}
u=\prod_{\alpha>0}x_{\alpha}(c_{\alpha}),
\end{equation}
其中 $c_{\alpha}\in F$，$x_{\alpha}: \Ga\to N$ 是 $\alpha$ 对应的根群。为说明 $X$ 包含在某个仿射 Schubert 簇中，只需对满足 $[u]\in X(k)$ 的任意 $u\in N(F)$，给出 \eqref{root gp} 中出现的 $c_{\alpha}$ 的赋值下界。可以把 $\Ad(u^{-1})\gamma\in\frb(F)$ 按根分解展开：
\begin{equation*}
\Ad(u^{-1})\gamma=\gamma+\sum_{\alpha>0}P_{\alpha}(\gamma; c)e_{\alpha},
\end{equation*}
其中 $e_{\alpha}\in\frg_{\alpha}$ 是每个根空间固定的基，而 $P_{\alpha}(\gamma;c)$ 是关于 $\{c_{\alpha}\}_{\alpha>0}$ 的 $F$-值多项式函数，并对 $\gamma$ 线性。对 $\alpha$ 的高度作归纳可知 $P_{\alpha}(\gamma;c)$ 具有如下形式
\begin{equation}\label{conj u gamma}
P_{\alpha}(\gamma;c)=\jiao{\alpha,\gamma}c_{\alpha}+\sum_{\beta<\alpha}\jiao{\beta,\gamma}P^{\beta}_{\alpha}(c),
\end{equation}
其中 $P^{\beta}_{\alpha}(c)$ 是只涉及 $\{c_{\alpha'}\}_{\alpha'<\alpha}$ 的多项式，且关于次数 $\alpha$ 齐次（这里约定 $\deg c_{\alpha'}:=\alpha'\in\xch(T)$）。

令 $n=\max_{\alpha>0}\{\val_{F}\jiao{\alpha,\gamma}\}$。此值有限，因为 $\gamma$ 正则半单。若 $[u]\in X(k)$，即 $\Ad(u^{-1})\gamma\in\frg(\calO_{F})$，则对 $\alpha$ 的高度作归纳可知
\begin{equation*}
\val_{F}(c_{\alpha})\geq -(2\textup{ht}(\alpha)-1)n,
\end{equation*}
这就给出了所需的下界，并说明 $X$ 落在某个仿射 Schubert 簇中。

\subsubsection{一般情形下 (1) 的证明} 一般情形下，我们给出一个比 \cite{KL} 中原本的论证更简化的论证，思路相同。作基变换到 $F'=k\lr{t^{1/m}}$，在其上 $\gamma$ 可以共轭进分裂环面。令 $\Gr'_{G}$ 为用域 $F'$ 代替 $F$ 定义的 $G$ 的仿射 Grassmann 簇（于是 $\Gr'_{G}(k)=G(F')/G(\calO'_{F})$），并令 $\sX'_{\gamma}\subset\Gr'_{G}$ 为相应的仿射 Springer 纤维。则 $\Gr'_{G}$ 与 $\sX'_{\gamma}$ 都带有 $\Gamma:=\Gal(F'/F)\cong\mu_{m}$ 的作用，该作用由 $\Gamma$ 在 $F'$ 上的作用诱导，并且有闭嵌入 $\sX_{\gamma}\incl (\sX'_{\gamma})^{\Gamma}$。

令 $\L'_{\gamma}=\xcoch(G_{\gamma}\otimes F')\incl G_{\gamma}(F')$ 为用域 $F'$ 构造的格。$\Gamma$ 作用在 $\xcoch(G_{\gamma}\otimes F')$ 上，不动点为 $\xcoch(G_{\gamma})\cong \L_{\gamma}$。此作用诱导 $\Gamma$ 在 $\L'_{\gamma}$ 上的作用，但它不保持嵌入 $\xcoch(G_{\gamma}\otimes F')\incl G_{\gamma}(F')$。不动点 $(\L'_{\gamma})^{\Gamma}=\xcoch(G_{\gamma}\otimes F')^{\Gamma}$ 未必落在 $G_{\gamma}(F)$ 中，但总以有限指数包含 $\L_{\gamma}$。

由 \S\ref{sss:split case} 中分裂情形的证明，有来自 Iwasawa 分解的有限型局部闭子概形 $X'\subset\sX'_{\gamma}$，使 $\sX'_{\gamma}$ 可以分解为{\em 不交并} $\sqcup_{\ell\in\L'_{\gamma}}\ell\cdot X'$（不是作为概形，而是作为可构造集）。把 $\ell\in \L'_{\gamma}$ 与 $\xcoch(G_{\gamma}\otimes F')$ 中的元素等同，则对 $\sigma\in\Gamma$，$\sigma(\ell)$ 有意义。容易验证：对 $\ell\in\L'_{\gamma}$ 有 $\sigma(\ell\cdot X')=\sigma(\ell)\cdot X'$，尽管 $\sigma$ 不保持嵌入 $\L'_{\gamma}\to G_{\gamma}(F')$。因此 $(\sX'_{\gamma})^{\Gamma}\subset (\L'_{\gamma})^{\Gamma}\cdot X'$。取有限陪集空间 $(\L'_{\gamma})^{\Gamma}/\L_{\gamma}$ 的代表元集 $C$，便知 $(\sX'_{\gamma})^{\Gamma}$ 包含在 $\L_{\gamma}\cdot (C\cdot X')$ 中。由于 $C\cdot X'$ 有限型，$Z'=\ov{C\cdot X'}\cap (\sX'_{\gamma})^{\Gamma}$ 是 $\sX'_{\gamma}$ 的射影子概形，其 $\L_{\gamma}$-平移覆盖 $(\sX'_{\gamma})^{\Gamma}$。最后，$\sX_{\gamma}$ 的射影子概形 $Z=Z'\cap \sX_{\gamma}$ 满足 (1) 的要求。定理 \ref{th:ft quot} 证毕。\qed

\subsubsection{归约到 Levi} \S\ref{sss:split case} 中分裂情形下定理 \ref{th:ft quot} 的证明给出了更多信息。在 Iwasawa 分解中，对 $\l\in\xcoch(T)$，令 $S_{\l}\subset \Gr_{G}$ 为 $\l(t)$ 的 $LN$-轨道。它称为{\em 半无限轨道}，因为它的维数无穷，在 $\Gr_{G}$ 中的余维数也无穷。令 $C_{\l}:=\sX_{\gamma}\cap S_{\l}$，则按 \S\ref{sss:split case} 的记号 $C_{\l}=\l(t)\cdot X$。公式 \eqref{conj u gamma} 表明：$C_{\l}\neq\varnothing$（等价地 $X\neq\varnothing$，或 $\sX_{\gamma}\neq\varnothing$）当且仅当对一切根 $\alpha$ 有 $\jiao{\alpha,\gamma}\in\calO_{F}$；而且此时 $C_{\l}$ 同构于一个维数为
\begin{equation*}
\dim C_{\l}=\sum_{\alpha>0}\val_{F}(\jiao{\alpha,\gamma})=\frac{1}{2}\val_{F}\Delta(\gamma)
\end{equation*}
的几乎仿射空间（即迭代的 $\AA^{1}$-丛）。这里 $\Delta(\gamma)$ 是 $\gamma$ 在 $\frg(F)/\frt(F)$ 上伴随作用之行列式。因此 $\sX_{\gamma}$ 可以分解为同维数的几乎仿射空间之并，以 $\xcoch(T)$ 为指标。但这个分解不是分层：$C_{\l}$ 的闭包 $\ov{C}_{\l}$ 会与其他 $C_{\l'}$ 相交，却肯定不是若干 $C_{\l'}$ 之并。

分裂情形的分解 $\sX_{\gamma}=\sqcup C_{\l}$ 有推广。设 $P$ 是 $G$ 的抛物子群，幂幺根为 $N_{P}$，Levi 子群为 $M_{P}$。令 $\mathfrak{m}_{P}=\Lie M_{P}$，并设 $\gamma\in\mathfrak{m}_{P}(F)$，且作为 $\frg(F)$ 中的元素正则半单。利用广义 Iwasawa 分解 $G(F)=N_{P}(F)M_{P}(F)G(\calO_{F})$，有定义良好的映射 $\Gr_{G}(k)\to \Gr_{M_{P}}(k)$，它把 $nm G(\calO_{F})$（$n\in N_{P}(F)$，$m\in M_{P}(F)$）送到 $m M_{P}(\calO_{F})$。然而此映射不给出归纳概形的映射。尽管如此，此映射的纤维具有无穷维仿射空间的自然结构。把此映射限制到 $\sX_{\gamma}$，得到
\begin{equation*}
\tau: \sX_{\gamma}(k)\to \sX^{M_{P}}_{\gamma}(k),
\end{equation*}
其中 $\sX^{M_{P}}_{\gamma}\subset \Gr_{M_{P}}$ 是 $\gamma$ 与群 $M_{P}$ 的仿射 Springer 纤维。$\tau$ 的纤维若非空，则是维数为 $\frac{1}{2}\val_{F}\Delta^{G}_{M_{P}}(\gamma)$ 的几乎仿射空间，其中
\begin{equation}\label{Delta G L}
\Delta^{G}_{M_{P}}(\gamma):=\det(\ad(\gamma)|\frg(F)/\mathfrak{m}_{P}(F)).
\end{equation}
设 $T\subset M_{P}$，则 $\Gr_{M_{P}}$ 的连通分支以 $\xcoch(T)/R^{\vee}_{M_{P}}$ 为指标，其中 $R^{\vee}_{M_{P}}$ 是 $M_{P}$ 的余根格。若把 $\sX^{M_{P}}_{\gamma}$ 按 $\xcoch(T)/R^{\vee}_{M_{P}}$ 分解为连通分支 $\sX^{M_{P}}_{\gamma}(\l)$ 并取它们在 $\tau$ 下的原像，就得到局部闭子归纳概形之分解 $\sX_{\gamma}=\sqcup\sX_{\gamma,\l}$，指标同样是 $\l\in\xcoch(T)/R^{\vee}_{M_{P}}$。可以证明：$\sX_{\gamma,\l}$ 若非空，则是 $\sX^{M_{P}}_{\gamma}(\l)$ 上的几乎仿射空间丛，纤维维数为 $\frac{1}{2}\val_{F}\Delta^{G}_{M_{P}}(\gamma)$。

\subsubsection{连通性与等维性}\label{sss:equidim} 当 $G$ 单连通时，$\Fl$ 连通，此时仿射 Springer 纤维 $\sY_{\gamma}$ 也连通。参见 \cite[\S4, 引理 2]{KL}。于是当 $G$ 单连通时，对一切 parahoric 子群 $\bP$，$\sX_{\bP,\gamma}$ 都连通，因为自然投影 $\sY_{\gamma}\to \sX_{\bP,\gamma}$ 满射。

\cite{KL} 中还证明了 $\sY_{\gamma}$ 等维。那里的论证类似 Spaltenstein 在 \cite{Spa77} 中对 Springer 纤维的连通性与等维性的证明。

\subsubsection{维数公式}
由定理 \ref{th:ft quot}，$\sX_{\gamma}$ 的维数有定义，且等于代数空间 $L_{\gamma}\bs\sX_{\gamma}$ 的维数。为陈述 $\dim\sX_{\gamma}$ 的公式，还需要一些记号。

考虑伴随作用 $\ad(\gamma):\frg(F)\to \frg(F)$。此映射的核是 $\frg_{\gamma}(F)$，而诱导的自同态 $\ov{\ad(\gamma)}$ 在 $\frg(F)/\frg_{\gamma}(F)$ 上可逆。令 $\Delta(\gamma)\in F^{\times}$ 为 $\ov{\ad(\gamma)}$ 的行列式。这与先前在 $\gamma$ 落在分裂环面 $\frt(F)$ 中时给出的 $\Delta(\gamma)$ 的定义一致\zy{$\Delta(\gamma)$ 与 \S4.2 中判别式 $\Delta(P_{a})$ 用同一字母而含义不同：前者是 $\ad(\gamma)$ 在 $\frg(F)/\frg_{\gamma}(F)$ 上的行列式，后者是多项式 $P_{a}$ 的判别式。来源：原文 \S2.5 该段与 \S4.2 文脉。}。

另一方面，回忆 $\xcoch(G_{\gamma})$ 是 $G_{\gamma}$ 的 $F$-有理余权群，它也是 $G_{\gamma}$ 的最大 $F$-分裂子环面的秩。令
\begin{equation*}
c(\gamma)=r-\rk_{\ZZ}\xcoch(G_{\gamma}).
\end{equation*}
则 $c(\gamma)$ 也是 $G_{\gamma}$ 的最大 $F$-各向异性子环面的秩。

\begin{theorem}[Bezrukavnikov {\cite{Bez}}，Kazhdan--Lusztig 猜想 \cite{KL}]\label{th:aff dim} 设 $\gamma\in \frg(F)$ 是正则半单元。则
\begin{equation}\label{aff Spr dim}
\dim\sX_{\gamma}=\frac{1}{2}(\val_{F}\Delta(\gamma)-c(\gamma)).
\end{equation}
\end{theorem}

\subsubsection{证明梗概}\label{dim formula discussion} 证明中起关键作用的是 $\sX_{\gamma}$ 的{\em 正则点}概念。有赋值映射 $\ev: \sX_{\gamma}\to [\frg/G]$，它把 $[g]\in\sX_{\gamma}\subset \Gr_{G}$ 送到 $\Ad(g^{-1})\gamma$ 模 $t$ 的约化，后者在 $G$ 的伴随作用下定义良好。称 $[g]\in \sX_{\gamma}$ 是{\em 正则点}，若 $\ev([g])$ 落在 $[\frg/G]$ 的开子叠 $[\frg^{\reg}/G]$ 中。令 $\sX^{\reg}_{\gamma}\subset\sX_{\gamma}$ 为 $\sX_{\gamma}$ 中由正则点构成的开子归纳概形。可以证明 $\sX^{\reg}_{\gamma}$ 非空。把 $\sX^{\reg}_{\gamma}$ 在 $\sY_{\gamma}$ 中的原像记为 $\sY^{\reg}_{\gamma}$，则投影 $\sY^{\reg}_{\gamma}\to\sX^{\reg}_{\gamma}$ 是同构。由于如 \S\ref{sss:equidim} 所述 $\sY^{\reg}_{\gamma}$ 等维，有 $\dim\sX^{\reg}_{\gamma}=\dim\sY^{\reg}_{\gamma}=\dim\sY_{\gamma}$。当然 $\dim\sX_{\gamma}\leq\dim \sY_{\gamma}$，因此必有 $\dim \sX^{\reg}_{\gamma}=\dim \sX_{\gamma}=\dim \sY_{\gamma}$。于是只需计算 $\sX^{\reg}_{\gamma}$ 的维数。

回忆在 \S\ref{sss:sym ASF} 中定义了局部 Picard 群 $P_{a}$。$P_{a}$ 在 $\sX_{\gamma}$ 上的作用保持 $\sX^{\reg}_{\gamma}$，事实上 $\sX^{\reg}_{\gamma}$ 是 $P_{a}$-主齐性空间。因此只需计算 $P_{a}=LG_{\gamma}/L^{+}J_{a}$ 的维数。

考虑投影 $\chi_{\frt}: \frt\to\frc=\frt\sslash W$。沿 $a=\chi(\gamma):\Spec\calO_{F}\to\frc$ 拉回，得到有限态射 $\chi_{a}: \chi^{-1}_{\frt}(a)=\Spec A\to \Spec\calO_{F}$，其中 $A$ 是有限平坦的 $\calO_{F}$-代数。群概形 $J_{a}$ 与 $J'_{a}=(\Res^{A}_{\calO_{F}}(T\otimes_{k}A))^{W}$ 之间有密切关系，这里 Weyl 群 $W$ 对角地作用在 $T$ 与 $A$ 上。事实上 $J_{a}$ 与 $J'_{a}$ 在特殊纤维中只差连通分支。特别地，$\dim P_{a}=\dim P'_{a}$，其中 $P'_{a}=LG_{\gamma}/L^{+}J'_{a}$。不难看出
\begin{equation*}
\dim P'_{a}=\dim (\frt\otimes \tilA)^{W}-\dim (\frt\otimes A)^{W}=\dim (\frt\otimes (\tilA/A))^{W},
\end{equation*}
其中 $\tilA$ 是 $A$ 的正规化。由此即得维数公式 \eqref{aff Spr dim}。\qed

\begin{remark} 利用仿射 Springer 纤维与 Hitchin 纤维的联系，Ng\^o \cite[推论 4.16.2]{NgoFL} 证明了 $\sX^{\reg}_{\gamma}$ 实际上在 $\sX_{\gamma}$ 中稠密。特别地，$\sX_{\gamma}$ 的每个不可约分支都是有理簇。然而对更一般的仿射部分旗簇中的仿射 Springer 纤维，这一有理性不再成立，如 \S\ref{sss:BK example} 中 Bernstein--Kazhdan 的例子所示。
\end{remark}

\subsubsection{纯性} Goresky、Kottwitz 与 MacPherson \cite{GKM} 猜想仿射 Springer 纤维的上同调应当是纯的（若 $k=\overline{\FF}_{p}$ 指 Frobenius 权的意义，若 $k=\CC$ 指 Hodge 结构的含义）。仿射 Springer 纤维的纯性将使 \cite{GKM} 的作者们能用等变上同调中的局部化技巧证明非分歧元素的基本引理。这一纯性猜想在一般情形仍悬而未决。在 \cite{GKMPurity} 中，一类称为{\em 等值}的仿射 Springer 纤维被证明是纯的。

\subsubsection{扰动下的不变性} 设 $a,a'\in\frc(\calO_{F})\cap\frc^{\rs}(F)$。若 $a$ 与 $a'$ 在映射 $\frc(\calO_{F})\to\frc(\calO_{F}/t^{N})$ 下的像相同，则称 $a\equiv a'\mod t^{N}$。\cite[命题 3.5.1]{NgoFL} 证明了：对固定的 $a\in\frc(\calO_{F})\cap\frc^{\rs}(F)$，存在 $N$（依赖于 $a$），使只要 $a'\equiv a\mod t^{N}$，就有与作用相容的同构
\begin{equation*}
\tX_{a}\cong\tX_{a'} \textup{ 与 } P_{a}\cong P_{a'}.
\end{equation*}
因此可以说：在 $\frc(\calO_{F})$ 的 $t$-进拓扑下，$\sX_{a}$ 局部常值地依赖于 $a$。

例如取 $G=\GL_{n}$，并设 $a\in\frc(\calO_{F})$ 对应特征多项式 $P(x)=x^{n}+a_{1}x^{n-1}+\cdots+a_{n}$，其根都在 $\calO_{F}$ 中且模 $t$ 互不相同。则对任意 $a'\equiv a\mod t$，$a'$ 的特征多项式也模 $t$ 有互不相同的根。此时 $\tX_{a}$ 与 $\tX_{a'}$ 都是 $P_{a}\cong P_{a'}\cong \Gr_{T}$-主齐性空间。
